\BourbakiExerciseGroup

\begin{enumerate}
\def\labelenumi{\arabic{enumi})}
\item
  设 X, Y 是两个拓扑空间，A 是 X 的子集，B 是 Y 的子集。证明 \(\mathrm{Fr}(\mathbf{A} \times \mathbf{B}) = (\mathrm{Fr}(\mathbf{A}) \times \overline{\mathbf{B}}) \cup (\overline{\mathbf{A}} \times \mathrm{Fr}(\mathbf{B}))\)。
\item
  设 \(\mathbf{X}\), \(\mathbf{Y}\) 是两个拓扑空间，A 是 \(\mathbf{X} \times \mathbf{Y}\) 的闭子集。若 \(\mathfrak{B}(x)\) 是点 \(x \in \mathbf{X}\) 的所有邻域所成的集，证明
\end{enumerate}

\[
\bigcap_ {\mathbf {V} \in \mathfrak {B} (x)} \overline {{\mathrm{A} (\mathbf {V})}} = \mathrm{A} (x).
\]

给出积 \(X \times Y\) 的一个开子集 B 的例子，使得

\[
\bigcap_ {\mathbf {V} \in \mathfrak {B} (x)} \overline {{\mathbf {B} (\mathbf {V})}}
\]

异于 \(\mathbf{B}_{(x)}\)。

\begin{enumerate}
\def\labelenumi{\arabic{enumi})}
\setcounter{enumi}{2}
\tightlist
\item
  设 \((\mathbf{X}_t)_{i\in I}\) 是拓扑空间的一个无限族，使得每个 \(\mathbf{X}_t\) 含有（至少）两个不同的点 \(a_i, b_i\)，且存在 \(b_i\) 的一个邻域，它不含 \(a_i\)。
\end{enumerate}

\begin{enumerate}
\def\labelenumi{\alph{enumi})}
\tightlist
\item
  对每个指标 \(x \in I\)，设 \(c_{x}\) 是积空间
\end{enumerate}

\[
\mathbf {X} = \prod_ {k + 1} \mathbf {X} _ {i}
\]

中满足 \(\mathrm{pr}_x c_x = b_x\) 且 \(\mathrm{pr}_t c_x = a_t\)（只要 \(i \neq x\)）的点。证明由诸点 \(c_x\) 组成的集 C 的每一点都是孤立的。

\begin{enumerate}
\def\labelenumi{\alph{enumi})}
\setcounter{enumi}{1}
\item
  由 a) 推出：X 的拓扑具有可数基，当且仅当 I 可数且每个 \(X_{i}\) 的拓扑具有可数基（用 §1 的习题 7）。
\item
  证明：若指标集 I 不可数，则 X 的点 \(b = (b_{i})\) 没有可数的基本邻域系。
\end{enumerate}

\begin{enumerate}
\def\labelenumi{\arabic{enumi})}
\setcounter{enumi}{3}
\tightlist
\item
  设 K 是由两个数 0, 1 组成的离散空间；设 A 是无限集，X 是积空间 \(\mathbf{K}^{\mathbf{A}}\)。设
\end{enumerate}

\[
\mathbf {V} = \prod_ {\alpha \in \mathbf {A}} \mathbf {V} _ {\alpha}
\]

是 X 中的一个初等集；若 h 是这样的指标 \(\alpha\)（\(V_{\alpha} \neq K\)）的个数，令 \(\mu(V) = 2^{-h}\)。

\begin{enumerate}
\def\labelenumi{\alph{enumi})}
\tightlist
\item
  证明：若 \(\mathbf{U}_1, \ldots, \mathbf{U}_n\) 是两两不相交的初等集，则
\end{enumerate}

\[
\sum_ {k = 1} ^ {n} \mu (\mathrm{U} _ {k}) \leqslant \mathrm{I}.
\]

{[}把每个 \(U_{k}\) 写成 \(W_{k} \times K^{B}\) 的形式，其中 B（与 k 无关）是 A 的一个有限子集的余集，而 \(W_{k}\) 是有限集；数一下 \(W_{k}\) 中点的个数{]}。

\begin{enumerate}
\def\labelenumi{\alph{enumi})}
\setcounter{enumi}{1}
\item
  由 a) 推出 X 满足 §1 习题 7 的条件 \((\mathbf{D}_{\mathbf{IV}})\)。{[}若存在两两不相交的初等集的一个不可数族 \((\mathbf{U}_{\lambda})\)，则将存在一个整数 p \textgreater{} 0，使得 \(\mu(\mathbf{U}_{\lambda}) \geqslant 1/p\) 对无穷多个指标 \(\lambda\) 成立。{]}
\item
  证明：若 A 不可数，则 X 不满足 §1 习题 7 的条件 \((\mathbf{D}_{\mathrm{III}})\)（参见习题 3）。
\end{enumerate}

\begin{enumerate}
\def\labelenumi{\arabic{enumi})}
\setcounter{enumi}{4}
\tightlist
\item
  设 K 是由两个数 0, 1 组成的离散空间；设 A 是无限集，X 表示积空间 \(\mathbf{K}^{\mathfrak{P}(\mathbf{A})}\)。a) 设 J 是 A 的任一有限子集；设 J 有 p 个元素。对 J 的每个子集 H，设 \(M_{H}\) 是 A 的满足 \(L \cap J = H\) 的子集 L 所成的集；诸 \(M_{H}\) 构成一个分拆 \(\tilde{\omega}_{J}\)（把 \(\mathfrak{P}(A)\) 分成 \(2^{p}\) 个子集）。设 \(F_{J}\) 是 X 中由这样的点 \((x_{L})_{L \subset A}\) 组成的子集：\(x_{L} = x_{M}\)（只要 L 与 M 属于分拆 \(\tilde{\omega}_{J}\) 的同一个集合）。\(F_{J}\) 是有限集，有 \(2^{2^{p}}\) 个元素。若 F 是当 J 取遍 A 的所有有限子集时诸 \(F_{J}\) 的并，证明 F 与 A 对等。
\end{enumerate}

\begin{enumerate}
\def\labelenumi{\alph{enumi})}
\setcounter{enumi}{1}
\tightlist
\item
  证明 F 在 X 中稠密。取 A = N，由此推出 X 满足 §1 习题 7 的条件 \((D_{\mathrm{II}})\)，但不满足条件 \((D_{\mathrm{III}})\)（习题 4）。
\end{enumerate}

\begin{enumerate}
\def\labelenumi{\arabic{enumi})}
\setcounter{enumi}{5}
\item
  设 X, Y, Z 是三个拓扑空间，A 是 \(X \times Y\) 的开子集，B 是 \(Y \times Z\) 的开子集。证明 \(B \circ A\) 是 \(X \times Z\) 的开子集。
\item
  设 X 是拓扑空间族 \((X_{\lambda})\) 的和，Y 是拓扑空间族 \((Y_{\mu})\) 的和。证明积空间 \(X \times Y\) 是诸 \(X_{\lambda} \times Y_{\mu}\) 的和。
\item
  设 \((\mathbf{G}_t)_{t \in \mathbf{I}}\) 是拓扑空间 \(\mathbf{X}\) 上一些等价关系的图象（《集合论》，第二章，§ 6，第 1 目）组成的族，且设 \(\mathbf{R}_t\) 是关系 \((x, y) \in \mathbf{G}_t\)；则
\end{enumerate}

\[
\mathbf {G} = \bigcap_ {\iota \in \mathbf {I}} \mathbf {G} _ {\iota}
\]

是 X 上一个等价关系 R 的图象。证明存在商空间 X/R 到积空间 \(\prod_{i}\left(X/R_{i}\right)\) 中的一个典范连续单射。

\begin{enumerate}
\def\labelenumi{\arabic{enumi})}
\setcounter{enumi}{8}
\tightlist
\item
  设 \((\mathbf{X}_i)_{i\in I}\) 是拓扑空间的一个无限族。研究积集合上的拓扑
\end{enumerate}

\[
\mathbf {X} = \prod_ {i \in I} \mathbf {X} _ {i} \text {   generated   by   subsets   of   the   form   } \prod_ {i \in I} \mathbf {U} _ {i},
\]

其中 \(U_{t}\) 是 \(X_{t}\) 的开子集（对每个 \(t \in I\)），且使 \(U_{t} \neq X_{t}\) 的指标 t 所成的集的基数小于 c，这里 c 是给定的无限基数。考察 §4 的诸命题在这个拓扑下变成什么。

\begin{enumerate}
\def\labelenumi{\arabic{enumi})}
\setcounter{enumi}{9}
\tightlist
\item
  \begin{enumerate}
  \def\labelenumii{\alph{enumii})}
  \tightlist
  \item
    设 \((\mathbf{X}_{\alpha}, f_{\alpha\beta})\) 是拓扑空间的一个逆系统，X 是逆极限。设 \(f_{\alpha}\) 是典范映射 \(X \to X_{\alpha}\)。对每个 \(\alpha\)，设 \(Y_{\alpha}\) 是 \(X_{\alpha}\) 的一个含有 \(f_{\alpha}(X)\) 的子空间，使得诸 \(Y_{\alpha}\) 构成诸 \(X_{\alpha}\) 的子空间的一个逆系统。证明 \(\varprojlim Y_{\alpha}\) 可典范地等同于 X。
  \end{enumerate}
\end{enumerate}

\begin{enumerate}
\def\labelenumi{\alph{enumi})}
\setcounter{enumi}{1}
\tightlist
\item
  设 \((\mathbf{X}_{\alpha}^{\prime}, f_{\alpha\beta}^{\prime})\) 是带有同一有向指标集的另一个拓扑空间逆系统。对每个 \(\alpha\)，设 \(u_{\alpha}: X_{\alpha} \to X_{\alpha}^{\prime}\) 是连续映射，使得诸 \(u_{\alpha}\) 构成映射的一个逆系统。诸 \(u_{\alpha}(\mathbf{X}_{\alpha})\) 构成诸 \(X_{\alpha}^{\prime}\) 的子空间的一个逆系统；证明：若 \(u = \varprojlim u_{\alpha}\) 且诸 \(f_{\alpha}\) 是满射，则 \(u(\mathbf{X})\) 在空间 \(\varprojlim u_{\alpha}(\mathbf{X}_{\alpha})\) 中稠密。若不再假设诸 \(f_{\alpha}\) 是满射，这个结果还成立吗？（参见《集合论》，第三章，§1，习题 32。）
\end{enumerate}

